% =====================================================================================
%  Sec. 10 -- CONCLUSION.  Replaces paper/conclusion.tex in full.
%
%  STRUCTURE, reordered 2026-09-23 and audited after writing.  The content is the same
%  as the previous version of this file; the ORDER is not.  Previously: bound first,
%  vacuity last, so the section read "here is my certificate and it does not work".
%  Now:
%    the OBSTRUCTION as the result  |  the two-sided bound as the instrument that
%    identifies the object obstructed  |  the sweep as the evidence, with the one
%    positive property and its own defeater  |  THREE limits, named and numbered in
%    prose  |  the open question  |  one closing sentence.
%  No promise of advantage anywhere.  The closing sentence must not be turned into one.
%
%  WHAT THIS REORDER MAY NOT DO, and did not: add a claim the manuscript does not
%  already make, turn a hedge into an assertion, or drop a number, a concession, a
%  priority attribution or a retraction.  Two concessions were moved IN from Secs. 3.1
%  and 5 rather than out: the explicit retraction of Theorem 1(iii) of versions 1-2,
%  and the ten-of-fourteen out-of-sample result for the two-parameter rule fitted in
%  1-w_S, with the count that moves between six and twelve.  Both make the section
%  longer and are kept for that reason.
%
%  EXTERNAL LABELS USED (all owned by other sections):
%    thm:leakage, sec:leakage  Sec. 3          [sec_3_body.tex]
%    sec:retraction            Sec. 3.1        [sec_3_1.tex]
%    thm:moments               Sec. 4          [sec_4_body.tex]
%    sec:frontier, sec:frontier:obstruction    [sec_5_body.tex]
%    sec:nogo, thm:witness, prop:leakinv       [sec_6_body.tex]
%    sec:prices                Sec. 7          [sec_7.tex]
%    sec:physics, tab:gaps     Sec. 8          [sec_8.tex]
%    sec:honest                Sec. 9          [sec_9.tex]
%  This file defines only: sec:conclusion.
%
%  \cite keys used here, all VERIFIED PRESENT in paper/bibliography.tex (2026-09-23):
%    benthien2007, nocera2016, kuehner1999, whitefeiguin2004, holzner2011, kpm2006,
%    paeckel2019, tdasci2024, whitedmrg1992, schollwock2011, tarabunga2026.
%  The set is unchanged by the reorder: none added, none dropped.
%
%  GREP NOTES for the mechanical audit:
%   (a) "decoupled"/"decoupling" appear nowhere in this file.
%   (b) The section opens on the obstruction and does not end on one: the three limits
%       are followed by the open question and by what the paper leaves behind.
%   (c) No number retired on 2026-09-18 appears here.  In particular this file does NOT
%       say that the shot budget grows with a larger exponent than the support (Sec. 7
%       measures the difference at 0.054 per site with a 95% interval containing zero),
%       quotes no minimum slack for the bound, and quotes the spin-gap coefficient only
%       as the approximate value of Sec. 8, never to three figures.
%   (d) The headline impossibility carries its scope inside the same sentence ("at every
%       size this work reports"), and the L=6 support count (296 of 300, which is the
%       one size whose answer depends on the threshold) is printed next to it.  Neither
%       may be dropped: without them the sentence claims a theorem the paper does not
%       prove.
% =====================================================================================

\section{Conclusion}
\label{sec:conclusion}

The result of this paper is negative, and it is a statement about the problem rather than about our
inequality: from $L=8$ on, in the determinant basis a device measures in, no leakage-type
certificate---ours or another---can be non-vacuous at low fraction. The reason
is geometric and it is measured. Zero leakage requires the subspace to contain the Krylov space
$\mathcal K(H,\phi)$, and the coordinate support of that Krylov space is $296$ of $300$ determinants
at $L=6$ and all $3920$ of $3920$ at $L=8$---so at $L=8$ the leakage is strictly positive for every
proper coordinate subspace---and, at an amplitude threshold of $10^{-12}$, the whole sector at
$L=8$, $10$, $12$ and $14$, the largest of dimension $10\,306\,296$, by two estimators that agree
wherever both ran (Sec.~\ref{sec:frontier:obstruction}). The obstruction is structural rather than
incidental, and a property of the basis the device measures in, not of the bound and not of the
ranking.

The two-sided bound is what identifies the object that obstruction is about. By
Theorem~\ref{thm:leakage} the error of a spectral function reconstructed on a subspace of sampled
configurations splits into the Born weight the shots failed to capture and the Hamiltonian coupling
that leaks across the boundary of that subspace at the resolution asked for; only the first of the
two is bought by drawing more bitstrings, so the error is indexed on the boundary and not on the
weight. The weight alone cannot carry the statement: no inequality indexed on the captured weight can
hold, by an analytic counterexample and by counting, and the weight-only Theorem~1(iii) of
versions~1 and~2 of this preprint is retracted (Sec.~\ref{sec:retraction}). What a Krylov-adequate
subspace does guarantee is moment exactness through order $2K+1$, with the proof the code realizes
(Theorem~\ref{thm:moments}). And nothing cheaper bounds that boundary \emph{usefully} from outside: one-body
fermionic magic is constant along an orbit on which the determinant support runs from $O(K)$ to
$2^{K}$ at bond dimension $2$, and constant on fibers on which the subspace-aware certificate still
separates (Theorem~\ref{thm:witness}, Proposition~\ref{prop:leakinv}), the general
invariant/basis-dependent division being due to Tarabunga \emph{et al.}~\cite{tarabunga2026}.

Measured subspace by subspace at five system sizes, that bound is vacuous at every published
fraction, at each of the five sizes evaluated, by factors of $1.66$ to $8.61$, converged in Krylov
depth at $L=12$---not an accident of the constants, but the obstruction in the numbers. The same
sweep yields the one positive property that survived four independent attacks: the crossing of the
trivial bound occurs at a true relative $L_1$ error between $1.2\times10^{-3}$ and
$7.3\times10^{-3}$, median $2.8\times10^{-3}$, over sixteen $(L,\eta)$ cells at $L=6$--$12$ and
$\eta=0.10$--$0.25\,t$---a drift, not scatter: ordered without exception in both arguments and
conservative in $L$---and it is computed without the answer (Sec.~\ref{sec:frontier}). Its own limit
is in the same table: a rule with two constants fitted in $1-w_S$ localizes that crossing about as
tightly out of sample and more tightly on ten of the fourteen train/test splits, a count that moves
between six and twelve as the constants are fitted, so we claim no superiority for the certificate as
a tracker of the error against the captured weight---out of sample it has none. The method is priced
on three axes rather than one, the third---shots---measured here for the first time
(Sec.~\ref{sec:prices}), and the primitive is validated against the one answer known exactly and
independently of it: from the same calculation on the same lattice, the charge scale converges to a
finite thermodynamic value while the spin scale closes as $\approx2.1/L$, with an independent
effective-Heisenberg control (Table~\ref{tab:gaps}).

Three limits bound all of it, and we state them as flatly as the results. \emph{First, the
certificate does not certify the reconstructions this paper reports.} They are, at present,
uncertified---vacuous at every published fraction and at each of the five sizes evaluated,
above---and neither a sharper proof nor a rearrangement of the subspace repairs that.
\emph{Second, the largest momentum-resolved spectrum is $L=12$, by exact diagonalization.} The three
momentum-resolved responses reported here were computed classically by dynamical DMRG in
2007~\cite{benthien2007}---at $120$, $100$ and $90$ sites in the charge, spin and single-particle
channels respectively---and at $L=48$ in the Hubbard chain, $L=64$ in the Heisenberg chain, with
correction vectors~\cite{nocera2016}; the wider classical toolbox is broader
still---correction-vector and time-dependent DMRG, Chebyshev matrix-product states and
kernel-polynomial expansions~\cite{kuehner1999,whitefeiguin2004,holzner2011,kpm2006,paeckel2019}, and
the real-time propagation of selected configuration-interaction wavefunctions~\cite{tdasci2024}; one
dimension is where tensor networks, since White's formulation~\cite{whitedmrg1992}, are
near-optimal~\cite{schollwock2011}, and nothing in this work competes with that standard; the
vocabulary of the thermodynamic limit is accordingly absent from it (Sec.~\ref{sec:physics}).
\emph{Third, the certified regime costs more shots than the method saves.} Moving from the published
fraction to the non-vacuous band multiplies the shots per snapshot needed for the marginal
configuration by $7.5$ to $10$ at $L=8$ and by $20$ to $35$ at $L=12$, where that configuration
carries a Born probability of $1$ to $2\times10^{-9}$: a regime a device can certify is, on these
systems, one in which diagonalizing the whole sector is already the cheaper thing to do.

Two questions follow, and only one of them is a matter of more computing. Whether a certified regime
exists at low sector fraction is now a well-posed measurement rather than an expectation, and the
calibration above is the instrument for making it; whether it exists at all turns on a prior question
this work can pose but not answer---whether any single-particle basis leaves the coordinate support
of $\mathcal K(H,\phi)$ short of the entire sector, since it is that support, and not the quality of
the sampler, that pins the threshold measured here.

What this paper leaves behind is therefore narrow and, we think, useful: an obstruction measured
rather than conjectured, binding a class of certificates and not only ours; a two-sided bound that
names the boundary as the object the error is indexed on and can be evaluated without the answer; and
a map of where that bound is informative, its negative half stated in the same voice as its positive
one.

% ---------------------------------------------------------------------------------
%  ACKNOWLEDGMENTS -- added 2026-09-19 at the end of the last body section, which is
%  where REVTeX places this block: main.tex inputs sec_10.tex and then \clearpage
%  \appendix, so the block falls after the Conclusion and before Appendix A.
%
%  WHAT PRA ACTUALLY REQUIRES, checked rather than assumed (2026-09-19):
%    * Data Availability Statement: REQUIRED.  "All published articles must include a
%      Data Availability Statement (DAS)" -- journals.aps.org/pra/authors, and
%      journals.aps.org/authors/data-availability-statements.  Written as Sec. 9.6.
%    * Conflict of interest: NOT required inside the article.  APS requires authors to
%      "alert editors to any potential conflict of interest such as sources of funding,
%      condition of employment, and any limitations on disclosure of data, code, or
%      experimental details", and only ENCOURAGES a statement in the paper itself
%      ("authors are encouraged to declare any conflict of interest within the paper
%      itself using a Conflict of Interest statement") -- journals.aps.org/authors/
%      editorial-policies.  So none is printed here.  If Nicolas wants one, it goes in
%      this block; the employment affiliation is the item an editor would expect named.
%    * Funding: nothing is claimed here, in either direction.  This block asserts only
%      what can be verified from the deposit, i.e. that the device runs happened.
%  The IBM sentence is IBM Quantum's standard acknowledgement and disclaimer.
% ---------------------------------------------------------------------------------
\begin{acknowledgments}
We acknowledge the use of IBM~Quantum services for this work; the two device runs reported in
Sec.~\ref{sec:scope:hw} were executed on \texttt{ibm\_fez} and \texttt{ibm\_marrakesh}. The views
expressed are those of the author and do not reflect the official policy or position of IBM or the
IBM~Quantum team.
\end{acknowledgments}
